\section{Hạn chế và mối đe doạ tính hợp lệ}\label{sec:limitations}

\paragraph{Tính hợp lệ nội tại.}
\begin{itemize}
  \item \textbf{Giám khảo LLM.} Mọi số accuracy đến từ một giám khảo (\texttt{gemini-flash-lite-latest}), với một lượt chấm cho mỗi lượt sinh.
  Rubric gốc khớp với một người chấm ở \humanAgreeOrig{} câu Null, và \judgeLenient{} câu được chấm đúng vẫn chứa chi tiết mâu thuẫn với ngữ cảnh khi rà tay.
  Vì vậy accuracy có thể cao hơn độ trung thành thực tế. Nếu độ dễ dãi này khác nhau giữa các cấu hình, các hiệu số có thể bị lệch theo hướng chưa biết.
  Việc kiểm với người mới làm trên nhóm Null, và chỉ có một người chấm.
  \item \textbf{Mốc MiniRAG chỉ có một lượt sinh.} Các phép so với MiniRAG (X1--X3) không nằm trong đăng ký trước, và phía MiniRAG không có lượt lặp.
  Vì hiệu số lớn gấp \ratioXCT{} lần ngưỡng nhiễu sinh $T$, chúng tôi cho rằng hướng của kết luận đứng vững, nhưng độ lớn chính xác cần thêm lượt sinh để xác nhận.
  \item \textbf{Lệch đăng ký.} RRF(vectơ, BM25) vào vòng xác nhận sau khi cách đo cổng thời gian được sửa, và việc sửa diễn ra sau khi cấu hình này trượt cổng (\cref{sec:prereg}).
  H2 trượt cổng an toàn của nhóm Null và được giữ ở mức biên.
  \item \textbf{Tập dev được nhìn nhiều lần.} Toàn bộ chẩn đoán, sàng lọc và lựa chọn biến thể làm trên tập dev, nên \cref{tab:retrieval} chỉ là dự báo.
  Tập ngoài dev chỉ dùng cho vòng xác nhận, nhưng ba lượt RRF(đồ thị, vectơ) và một lượt vectơ thuần trên tập này đã được xem trước khi đăng ký BM25.
  \item \textbf{Nhiều cấu hình đã thử.} Chúng tôi đã thử 23 cấu hình hoặc hướng (\cref{app:tried}). Hiệu chỉnh Holm chỉ áp dụng trong từng họ phép so, không trên toàn dự án.
\end{itemize}

\paragraph{Tính hợp lệ ngoại tại.}
\begin{itemize}
  \item \textbf{Một bộ dữ liệu, một ngôn ngữ.} Mọi kết luận chỉ áp dụng cho LiHua-World, một kho hội thoại tiếng Anh tổng hợp do chính nhóm tác giả MiniRAG xây dựng.
  Câu hỏi của nó thường chứa tên riêng xuất hiện nguyên văn, một điều kiện rất thuận lợi cho BM25.
  Chúng tôi đã chuẩn bị 2WikiMultiHopQA và MultiHop-RAG nhưng chưa chạy \todo{chạy hoặc giữ nguyên câu này}.
  \item \textbf{Một mô hình sinh, một đồ thị.} Vòng xác nhận chỉ dùng Qwen2.5-3B, với đồ thị do chính mô hình này dựng (\graphIsolated\% nút cô lập).
  Một đồ thị dày hơn, ví dụ do mô hình lớn hơn dựng, có thể cho đường đi tốt hơn.
  \item \textbf{Nhóm câu nhỏ.} Multi (\nHeldMulti{} câu) và Null (\nHeldNull{} câu) quá nhỏ để kết luận riêng.
\end{itemize}

\paragraph{Tái lập.}
\begin{itemize}
  \item \textbf{Khác bài gốc.} Chúng tôi dùng giám khảo khác (Gemini thay cho GPT), đồ thị dựng lại sau khi sửa lỗi lập chỉ mục $O(N^2)$, ngân sách ngữ cảnh 4.000 token và bản sửa khớp loại đáp án.
  Vì vậy số của bài gốc chỉ là tham chiếu.
  \item \textbf{Đồ thị thưa hơn bản gốc.} Mã gốc vô tình trích xuất lặp nên gom thêm thực thể; bản sửa của chúng tôi bỏ việc lặp này.
  Điều này có thể làm MiniRAG tái lập yếu hơn bản gốc chạy đúng như công bố, và chúng tôi chưa đo được mức ảnh hưởng.
  \item \textbf{Nguồn ngẫu nhiên.} Lấy mẫu khi sinh đã được lặp ba lượt; giám khảo chỉ chấm một lượt cho mỗi lượt sinh; chỉ mục dựng bằng LLM chỉ dựng một lần.
\end{itemize}
